% Part: sets-functions-relations
% Chapter: relations
% Section: operations

\documentclass[../../../include/open-logic-section]{subfiles}

\begin{document}

\olfileid[fr]{sfr}{rel}{ops}
\olsection{Opérations sur les relations}

Il est souvent utile de modifier ou de combiner des relations.
Dans la \olref[sfr][rel][ord]{prop:stricttopartial}, nous avons considéré
leur \emph{réunion}, qui est simplement la réunion de deux relations
prises comme ensembles de couples. De même, dans la
\olref[sfr][rel][ord]{prop:partialtostrict}, nous avons considéré
leur différence ensembliste. Voici quelques autres opérations sur les relations.

\begin{defn}\ollabel{relationoperations}
Soient $R$ et $S$ des relations, et $A$ un ensemble quelconque.

La relation \emph{inverse} de $R$ est
$R^{-1} = \Setabs{\tuple{y, x}}{\tuple{x, y} \in R}$.

Le \emph{produit relatif} de $R$ et $S$ est
$(R \mid S) = \{\tuple{x, z} : \exists y(Rxy \land Syz)\}$.

La \emph{restriction} de $R$ à $A$ est
$\funrestrictionto{R}{A}= R \cap A^2$.

L’\emph{image} de $A$ par $R$ est
$\funimage{R}{A} = \{y : (\exists x \in A)Rxy\}$.
\end{defn}

\begin{ex}
Soit $S \subseteq \Int^2$ la relation de succession sur~$\Int$ :
$S = \Setabs{\tuple{x, y} \in \Int^2}{x + 1 = y}$.
Ainsi, $Sxy$ si et seulement si $x + 1 = y$.

$S^{-1}$ est la relation qui associe à chaque entier son prédécesseur,
soit $\Setabs{\tuple{x,y}\in\Int^2}{x -1 =y}$.

$S\mid S$ est $ \Setabs{\tuple{x,y}\in\Int^2}{x + 2 =y}$.

$\funrestrictionto{S}{\Nat}$ est la relation de succession sur~$\Nat$.

$\funimage{S}{\{1,2,3\}}$ est $\{2, 3, 4\}$.
\end{ex}

\begin{defn}[Clôture transitive]
Soit $R \subseteq A^2$ une relation binaire.\footnote{Précision de notation :
dans cette section, le signe plus en exposant désigne la clôture transitive.
Il n’a donc pas le sens local de clôture réflexive employé dans la
section sur les ordres. Les notations de l’original sont conservées.}

La \emph{clôture transitive} de~$R$ est
$R^+ = \bigcup_{0 < n \in \Nat} R^n$, où l’on définit par récurrence
$R^1 = R$ et $R^{n+1} = R^n \mid R$.

La \emph{clôture réflexive et transitive} de $R$ est
$R^* = R^+ \cup \Id{A}$.
\end{defn}

\begin{ex}
Reprenons la relation de succession $S \subseteq \Int^2$.
On a $S^2xy$ si et seulement si $x + 2 = y$, $S^3xy$ si et seulement si
$x + 3 = y$, etc. Ainsi, $S^+xy$ si et seulement si $x + n = y$
pour un certain $n \geq 1$. Autrement dit, $S^+xy$ si et seulement si
$x < y$, et $S^*xy$ si et seulement si $x \le y$.
\end{ex}

\begin{prob}
Montrez que la clôture transitive de $R$ est effectivement transitive.
\end{prob}

\end{document}
